\documentclass{article}
\usepackage{PRIMEarxiv} % Core package for the PrimeAI Template
\usepackage{amsmath, amssymb, amsthm, bm, dcolumn} % Add amsthm here for the proof environment
\usepackage[numbers,sort&compress]{natbib} % Natbib for citations
\usepackage{graphicx} % For high-quality images
\usepackage{multicol} % Multi-column figures/tables
\usepackage{hyperref} % Hyperlinks
\usepackage{listings} % Code listings
\usepackage{authblk} % For structured affiliations
% Define theorem style (optional)
\theoremstyle{plain}
\newtheorem{theorem}{Theorem}
\renewcommand{\qedsymbol}{}

% Custom commands for primes and qubit encoding
\newcommand{\primeQubit}[1]{|p_{#1}\rangle}
\newcommand{\primeGate}[1]{U_{p_{#1}}}

\usepackage{wrapfig}
\usepackage[pscoord]{eso-pic}
\usepackage[fulladjust]{marginnote}
\reversemarginpar

% Typesetting improvements without footnote patching
\usepackage[protrusion=true, expansion=true, tracking=false]{microtype}
\microtypecontext{spacing=nonfrench}

% Line numbers
\usepackage[right]{lineno}

% Text layout - adjust as needed
\raggedright
\setlength{\parindent}{0.5cm}
\textwidth 5.25in 
\textheight 8.75in

% Set double spacing
\usepackage{setspace} 
\doublespacing

% Adjust width for specific content
\usepackage{changepage}

% Adjust caption style
\usepackage[aboveskip=1pt,labelfont=bf,labelsep=period,singlelinecheck=off]{caption}
% Define custom claims environment
\newtheoremstyle{claimstyle}
  {3pt}   % Space above
  {3pt}   % Space below
  {\itshape}   % Body font
  {}      % Indent amount
  {\bfseries} % Theorem head font
  {.}     % Punctuation after theorem head
  { }     % Space after theorem head
  {\thmname{#1}\thmnumber{ #2}\thmnote{ (#3)}} % Theorem head spec

\theoremstyle{claimstyle}
\newtheorem{claim}{Claim}
% Remove brackets from references
\makeatletter
\renewcommand{\@biblabel}[1]{\quad#1.}
\makeatother

% Header, footer, and page numbers
\usepackage{lastpage,fancyhdr}
\pagestyle{fancy}
\fancyhf{}
\fancyhead[L]{Preprint - PrimeAI Enhanced Template}
\fancyfoot[C]{\scriptsize Multiplicity Theory © 2024 Ryan Van Gelder - Citizen Gardens \\ Licensed Under MIT and CC BY-NC-SA 4.0.}
\fancyfoot[R]{Page \thepage\ of \pageref{LastPage}}
\renewcommand{\footrule}{\hrule height 2pt \vspace{2mm}}

\begin{document}
\title{Prime-Indexed Recursive Tensor Mathematics (PIRTM) for Cardiovascular Disease Modeling}
\author{Ryan O. Van Gelder} 
\date{\today}



\maketitle

\begin{abstract}
Cardiovascular disease (CVD) arises from complex, nonlinear interactions among genetic, physiological, and environmental factors, challenging traditional models with unstable feedback loops and nonstationary dynamics. Prime-Indexed Recursive Tensor Mathematics (PIRTM) provides a stable, predictive framework for CVD modeling through recursive tensor updates, prime-weighted structures, and reinforcement learning (RL)-based optimization. This paper enhances PIRTM by integrating the Multiplicity Biological Model, incorporating Body Mass Index (BMI) and Electrocardiogram (EKG) data via prime-based bio-index encoding, recursive Bayesian learning, quantum-inspired feedback loops, and golden ratio tensor operators. The result is a robust, adaptive, and personalized approach to cardiovascular health monitoring and intervention design.
\end{abstract}

\section{Introduction}
Cardiovascular disease (CVD) progression involves multidimensional interactions that traditional models struggle to capture due to nonlinearity and instability. Prime-Indexed Recursive Tensor Mathematics (PIRTM) addresses these challenges by offering:
\begin{itemize}
    \item \textbf{Recursive Stability}: Mimics biological homeostasis through convergent tensor updates.
    \item \textbf{Hierarchical Prime-Indexed Weighting}: Distinguishes long-term (e.g., genetic) vs. short-term (e.g., environmental) influences.
    \item \textbf{Tensor-Based Predictive Dynamics}: Integrates high-dimensional data for forecasting.
\end{itemize}
By incorporating the Multiplicity Biological Model, we enhance PIRTM with:
\begin{itemize}
    \item \textbf{Real-Time Adaptability}: Uses BMI and EKG data for dynamic monitoring.
    \item \textbf{Quantum and Neuromorphic Insights}: Models cardiovascular entropy and coherence.
    \item \textbf{Personalized Predictions}: Employs Bayesian networks for probabilistic optimization.
\end{itemize}
This unified framework bridges mathematical innovation and biological intelligence for advanced CVD modeling.

\section{Mathematical Formulation}

\subsection{Tensor Representation of CVD State}
The CVD state is represented as a rank-$(m, n)$ tensor \( T_t^{(m, n)} \):
\begin{itemize}
    \item \( m \): Physiological factors (e.g., blood pressure (BP), LDL/HDL cholesterol, heart rate (HR), glucose levels, inflammation markers).
    \item \( n \): Temporal/external factors (e.g., genetic risk, environmental stress, medication).
\end{itemize}
Example:
\[
T_t^{(3, 2)} = \begin{bmatrix}
\text{BP} & \text{Genetic Risk} \\
\text{LDL} & \text{Smoking Status} \\
\text{HR} & \text{Diet Impact}
\end{bmatrix}.
\]

\subsection{Prime-Based Bio-Index Encoding}
Extend \( T_t^{(m, n)} \) with a bio-index tensor \( \mathcal{B}_t \) integrating BMI and EKG:
\begin{equation}
\mathcal{B}_t = \sum_{p_k \in P_N} T_{jk}^{(p_k)} p_k^{\alpha_k} (\xi(p_k) + \psi(p_k, t)),
\end{equation}
where:
\begin{itemize}
    \item \( \mathcal{B}_t \): Bio-index state at time \( t \).
    \item \( T_{jk}^{(p_k)} \): Prime-indexed health eigenvalues (e.g., BMI, HRV, BP).
    \item \( p_k^{\alpha_k} \): Prime-weighted scaling (\( \alpha_k < -1 \) for convergence).
    \item \( \xi(p_k) \): Systemic metabolic adaptation (e.g., BMI-driven metabolic rate).
    \item \( \psi(p_k, t) \): Self-referential cardiac dynamics (e.g., EKG variability).
\end{itemize}

\subsection{Recursive Tensor Evolution}
The enhanced PIRTM update equation is:
\begin{equation}
T_{t+1}^{(m, n)} = \sum_{p_i \in P_N} \Lambda_m \cdot p_i^{\alpha} \cdot T_t^{(m, n)} + F^{(m, n)} + \beta \mathcal{B}_t,
\end{equation}
where:
\begin{itemize}
    \item \( P_N = \{2, 3, 5, 7, 11\} \): Hierarchical primes.
    \item \( \Lambda_m \in (0, 1) \): Universal Multiplicity Constant.
    \item \( F^{(m, n)} \): External interventions (e.g., statins, exercise).
    \item \( \beta \): Weights bio-index contributions.
\end{itemize}

\subsection{Golden Ratio Tensor Operators}
Define HRV coherence with a golden ratio-inspired operator:
\begin{equation}
\Phi_T = \sum_{n} \frac{T_{ij}^{(n)}}{F_n},
\end{equation}
where \( F_n = \{1, 1, 2, 3, 5, \ldots\} \) approximates the Fibonacci sequence, linking cardiovascular stability to harmonic patterns.

\section{Stability and Homeostasis Analysis}
The system converges to:
\begin{equation}
\lim_{t \to \infty} T_t^{(m, n)} = T_{\infty}^{(m, n)} = \frac{F^{(m, n)} + \beta \mathcal{B}_{\infty}}{1 - k},
\end{equation}
where \( k = \sum_{p_i \in P_N} \Lambda_m p_i^{\alpha} \) and \( |k| < 1 \). \\
\textbf{Biological Interpretation:}
\begin{itemize}
    \item Healthy state: \( T_{\infty}^{(m, n)} \) reflects homeostasis (e.g., BP = 120/80 mmHg).
    \item Disease state: Deviations indicate CVD progression.
\end{itemize}

\subsection{Perturbation Analysis}
Model stressors as:
\begin{equation}
\epsilon_{t+1} = k \cdot \epsilon_t, \quad \epsilon_t = k^t \cdot \epsilon_0 \to 0 \text{ as } t \to \infty.
\end{equation}
Low \( |k| \) ensures rapid recovery; high \( |k| \) signals instability.

\section{Recursive Bayesian Learning}
Model probabilistic dependencies:
\begin{equation}
P(\text{BMI} | \text{EKG}, \text{Metabolic Rate}) = \frac{P(\text{EKG} | \text{BMI}) P(\text{BMI})}{P(\text{EKG})},
\end{equation}
updating \( T_t^{(m, n)} \) priors dynamically with real-time data.

\section{Quantum-Inspired Feedback Loops}
Introduce a feedback model:
\begin{equation}
M_{t+1} = f(M_t, R_t) + \alpha T(M_t),
\end{equation}
where \( M_t \) is the metabolic-cardiac state, \( R_t \) is real-time EKG/BMI input, and \( f(\cdot) \) is neuromorphic. \\
\textbf{Cardiovascular Entropy:}
\begin{equation}
H_{\text{prime}} = \sum_{p_i \in P_N} \Lambda_m e^{-\alpha p_i} \hat{H},
\end{equation}
where \( \hat{H} \) quantifies EKG-derived disorder.

\section{Reinforcement Learning for Personalization}
Enhance RL with:
\begin{itemize}
    \item \textbf{State}: \( T_t^{(m, n)} \).
    \item \textbf{Actions}: Adjust \( F^{(m, n)} \).
    \item \textbf{Reward}:
    \begin{equation}
    R_t = -\sum_{i,j} |T_{t,i,j} - T_{\text{healthy},i,j}| - \gamma H_{\text{prime}},
    \end{equation}
    penalizing entropy.
\end{itemize}
Policy weights update via:
\begin{equation}
W_{t+1} = \sum_{p_i \in P_N} \Lambda_m \cdot p_i^{\alpha} \cdot W_t + \eta \nabla L(W_t).
\end{equation}

\section{Computational Implementation}
\begin{verbatim}
import numpy as np
import matplotlib.pyplot as plt

# Parameters
N_primes = [2, 3, 5, 7, 11]
Lambda_m = 0.968
alpha = -1.2
k = Lambda_m * sum(p**alpha for p in N_primes)
beta = 0.1
steps = 50

# Initial state and interventions
T = np.array([140.0, 200.0, 75.0])  # BP, LDL, HR
F = np.array([-10.0, -20.0, -5.0])
B = np.array([30.0, 0.1])  # BMI, EKG variability

def bio_index(B, t):
    return B * (1 + 0.01 * np.sin(t))

def recursive_tensor_update(T, F, B, steps=steps):
    history = [T.copy()]
    for t in range(steps):
        B_t = bio_index(B, t)
        T = Lambda_m * sum(p**alpha * T for p in N_primes) + F + beta * B_t[0]
        history.append(T.copy())
    return np.array(history)

history = recursive_tensor_update(T, F, B)
fixed_point = (F + beta * B[0]) / (1 - k)

plt.figure(figsize=(10, 6))
labels = ['Blood Pressure (mmHg)', 'LDL Cholesterol (mg/dL)', 'Heart Rate (bpm)']
for i in range(3):
    plt.plot(history[:, i], label=labels[i])
    plt.axhline(y=fixed_point[i], color='gray', linestyle='--', alpha=0.5)
plt.title('CVD State Evolution with Enhanced PIRTM')
plt.xlabel('Time Step')
plt.ylabel('Value')
plt.legend()
plt.grid(True)
plt.show()

print("Initial State:", history[0])
print("Final State:", history[-1])
print("Fixed Point:", fixed_point)
\end{verbatim}

\section{Future Directions}
Explore:
\begin{itemize}
    \item Quantum biology modeling of oxidative stress.
    \item Topological analysis of vascular networks.
    \item Real-time wearable tech integration (e.g., ECG, BP monitors).
    \item Validation with clinical datasets (e.g., Framingham Heart Study).
\end{itemize}

\section{Conclusion}
This enhanced PIRTM framework, unified with the Multiplicity Biological Model, offers a mathematically rigorous, adaptive, and personalized approach to CVD modeling. By integrating prime-based bio-index encoding, Bayesian learning, quantum feedback, and golden ratio operators, it advances predictive accuracy and intervention design, paving the way for real-time health optimization.
\nocite{*}
\bibliographystyle{plain}
\bibliography{references}

\end{document}